\section{Pembahasan dan kesimpulan}\label{sec:discussion}
\subsection{Apa arti hasil bagi hipotesis awal}
Hasil belum memberi dasar bagi pernyataan bahwa depresiasi merupakan jaminan perbaikan neraca Indonesia. Agregat utama belum mendukung hubungan level yang kuat. Pada disaggregasi, terdapat pasar atau kelompok yang mengarah pada perbaikan dan terdapat yang mengarah pada pemburukan. Sebagian bukti lebih kuat menurut penyaringan yang dipakai, tetapi kesimpulannya tidak seragam antarspesifikasi.

Temuan linier komoditas total serta energi/tambang relevan bagi hipotesis awal. Arah negatif tetap terlihat pada beberapa pembanding, sehingga hubungan tersebut tidak hanya muncul pada satu tabel baseline. Namun NARDL belum memberikan konfirmasi setara, dan beberapa sensitivitas menolak diagnostik. Oleh karena itu, naskah tidak menyimpulkan bahwa karakter komoditas telah terbukti sebagai penyebab arah negatif tersebut.

Pangsa komoditas membantu menjelaskan mengapa agregasi barang layak dilakukan. Ia belum menghubungkan secara kausal harga internasional dengan koefisien kurs. Regresi tambahan BCOM pun belum memisahkan harga ekspor, harga impor, volume dan kondisi permintaan. Klaim mekanisme memerlukan pengukuran serta desain identifikasi yang berbeda dari sekadar menambahkan satu indeks kontrol.

\subsection{Batas penafsiran ekspor, neraca dan ketidakpekaan}
Outcome rasio neraca memberi informasi relatif. Penurunan impor yang lebih cepat daripada ekspor dapat menghasilkan perbaikan tanpa pertumbuhan ekspor. Sebaliknya, ekspor dapat meningkat tetapi impor naik lebih cepat, sehingga rasio memburuk. Menerjemahkan hasil menjadi pernyataan ekspor saja akan melewati informasi yang tersedia. Untuk hipotesis khusus ekspor, persamaan ekspor terpisah perlu dinilai dengan regressor dan identifikasi yang sesuai.

Ketidakpekaan terhadap kurs juga merupakan klaim yang lebih kuat daripada tidak signifikan. Interval lebar dapat memuat efek negatif maupun positif yang besar. Hasil seperti itu terutama menunjukkan ketidakpastian. Uji ekuivalensi akan memerlukan batas pengaruh yang dianggap kecil secara ekonomi dan justifikasi atas batas tersebut. Run ini belum melaksanakannya.

Heterogenitas antarbarang perlu diuji secara langsung apabila menjadi klaim utama. Membandingkan bintang pada dua regresi belum menguji kesamaan koefisien keduanya. Perbedaan window, korelasi residual, dan tumpang tindih kategori juga harus dipertimbangkan dalam desain uji. Bukti deskriptif dan matriks koefisien dalam PPT berfungsi memetakan masalah itu, tidak menggantikannya.

\subsection{Keterbatasan yang paling menentukan}
Pertama, revisi HS per sel belum diketahui. Konsensus lintas korespondensi menjaga agar kelas ambigu tidak dipaksakan, tetapi mengubah kode menjadi seri produk yang benar-benar sebanding sepanjang waktu masih memerlukan harmonisasi. Perubahan cakupan klasifikasi dapat berkaitan dengan komposisi kelompok dan memengaruhi rasio ekspor/impor.

Kedua, ketersediaan CPI, FX dan GDP menghasilkan sampel tidak seragam. India memiliki observasi lebih sedikit; celah AS dan Vietnam memotong window utama. Perbandingan lintas mitra karena itu mengandung perbedaan pengalaman ekonomi yang diamati. Proksi aktivitas agregat juga hanya merepresentasikan sepuluh mitra dengan bobot tetap.

Ketiga, nilai perdagangan belum dipisahkan ke harga dan volume serta belum disesuaikan musiman. Kejutan harga global dapat mengubah nilai dua arus meskipun volume tidak merespons. Pola musiman yang berbeda antara perdagangan dan GDP penyedia juga dapat berpengaruh. Diagnostik memberi petunjuk, tetapi tidak menentukan satu perbaikan yang pasti.

Keempat, kurs dapat merespons neraca dan kejutan makro yang sama. Lag membantu menggambarkan waktu penyesuaian, tetapi tidak otomatis membuat regressor eksogen. Tidak ada instrumen, kejutan kebijakan teridentifikasi atau rancangan kausal lain dalam run ini. Istilah hubungan kondisional karena itu lebih sesuai daripada efek kebijakan yang telah terisolasi.

Kelima, 516 estimasi menciptakan ruang besar untuk pencarian hasil. BH hanya diterapkan dalam keluarga utama tertentu dan belum menyelesaikan seleksi lag maupun semua perbandingan sensitivitas. Pelaporan lengkap merupakan langkah transparansi, bukan koreksi penuh seluruh ketidakpastian tersebut.

\subsection{Pengembangan penelitian yang paling dekat dengan pertanyaan}
Prioritas empiris berikutnya adalah menyelaraskan definisi dan sampel sebelum memperbanyak model. Harmonisasi HS dan dokumentasi cakupan akan memperkuat komposisi barang. Perbandingan pada window yang sama akan membantu memisahkan heterogenitas ekonomi dari perbedaan periode. Pemeriksaan musiman dan spesifikasi residual perlu mengikuti persoalan yang terlihat, dengan inferensi yang sesuai bila asumsi varians atau korelasi ditolak.

Bila fokus akhirnya adalah ekspor, data kuantitas atau ukuran harga yang sebanding akan membantu membedakan respons nilai dari volume. Model ekspor terpisah dapat menjawab pertanyaan yang tidak dapat dijawab rasio neraca. Estimasi tiap HS-6 juga memerlukan penanganan nol, perubahan kode, kecukupan sampel dan pengujian berganda; tidak semua seri produk dapat langsung dimasukkan ke model log yang sama. Hipotesis input impor tetap berada di luar run ini dan memerlukan sumber data lain.

\subsection{Kesimpulan}
Analisis perdagangan Indonesia menunjukkan bahwa satu koefisien agregat belum merangkum seluruh pengalaman pasar dan barang. Australia dan Malaysia memberikan dukungan lebih baik pada NARDL bilateral, dengan ketahanan signifikansi yang berbeda setelah BH. ARDL komoditas total serta energi/tambang mengarah pada pemburukan saat depresiasi, tetapi dukungan nonlinear dan diagnostik pembanding belum seragam. Pola J terlihat pada sebagian titik estimasi, sementara dukungan kedua fase pada interval lebih jarang.

Kesimpulan yang dapat dipertahankan adalah bahwa manfaat depresiasi terhadap neraca tidak terukur secara seragam dalam spesifikasi yang diperiksa. Penelitian ini belum membuktikan bahwa depresiasi selalu merugikan, bahwa komoditas pasti tidak sensitif, atau bahwa harga global telah menjadi penyebab yang teridentifikasi. Pembedaan antara hasil, kelayakan model dan penafsiran mekanisme merupakan dasar untuk melanjutkan penelitian dengan klaim yang lebih tepat.
